\medskip

\emph{Pour chacune des affirmations suivantes, indiquer si elle est vraie ou fausse. Chaque réponse doit être justifiée. Une réponse non justifiée ne rapporte aucun point.}

\medskip

L'espace est rapporté à un repère orthonormé \Oijk.

On considère la droite $(d)$ dont une représentation paramétrique est :


\[\left\{\begin{array}{l c l}
x&=&\phantom{-}3 - 2 t \\
y&=&-1 \\
z&=&\phantom{-}2 - 6 t
\end{array} \quad, \text { où } t \in \R\right.\]

On considère également les points suivants :

\begin{itemize}
\item A$(3~;~-3~;~-2)$
\item B$(5~;~-4~;~-1)$
\item C le point de la droite $(d)$ d'abscisse 2
\item H le projeté orthogonal du point B sur le plan $\mathcal{P}$ d'équation $x + 3z -7 = 0$
\end{itemize}

\medskip

\textbf{Affirmation 1}

La droite $(d)$ et l'axe des ordonnées sont deux droites non coplanaires.

\bigskip

\textbf{Affirmation 2}

Le plan passant par $A$ et orthogonal à la droite $(d)$ a pour équation cartésienne :

\[x + 3z + 3= 0\]

\bigskip

\textbf{Affirmation 3}

Une mesure, exprimée en radian, de l'angle géométrique $\widehat{\mathrm{BAC}}$ est $\dfrac{\pi}{6}$.

\bigskip

\textbf{Affirmation 4}

La distance BH est égale à $\dfrac{\sqrt{10}}{2}$.
